\section{Gap Amplification for $\QMA^+(k)$ up to Criticality and Almost-$\QMA(k)=\NEXP$}\label{sec:protocol}

In the previous section, we described a very strong version of swap test and product test, noting that our disentangler channel has a good synergy with the new tests to overcome the drawbacks in their original versions.
%
In this section, we put the tools in the context of quantum Merlin-Arthur games with unentangled provers,
establishing our main complexity results~\cref{thm:almost_gen_amp,thm:qma3}.

\subsection{Gap Amplification for $\QMA^+(k)$ up to Criticality}

The gap amplification for $\QMA^+(k)$ is much less straightforward than $\QMA(2)$.
%
Indeed, a full gap amplification would imply $\QMA(2) = \NEXP.$
%
To give our half gap amplification promised in~\cref{thm:qma3}, we start by showing how to
simulate a $\QMA^+(k)$ protocol $\calP$ given the following kinds of proofs:
\begin{enumerate}
    \item\label{enu:nonnegative-state} one nonnegative-amplitudes proof $\ket{\psi}$;
    \item\label{enu:many-general-states} abundant equal copies of an arbitrary proofs over reals $\ket{\phi}$.
\end{enumerate}   
Note we are relaxing $k$ nonnegative-amplitudes proofs in a $\QMA^+(k)$ protocol with only one nonnegative-amplitudes
proof and general-amplitudes states.
%
The motivation is, roughly, to remove as many nonnegative-amplitudes proofs in a $\QMA^+(k)$ protocol as possible, so we get
closer to a general $\QMA(k)$ protocol. 

We will check whether  $\ket{\phi}^{\otimes k}$
is close to $\ket{\psi}$. Either they are close and then we can use the many copies of $\ket{\phi}^{\otimes k}$
to simulate $\calP$, or else they are far apart and an application of the super product test can detect this
condition. A description of this simulation procedure is given in~\cref{proc:sym_protocol}, which we denote
as the symmetric simulator (since it assumes many equal copies of $\ket{\phi}$).

\begin{algorithm}[H]{SymSimulator}{($\calP$, $\ket{\psi} = \sum_i \beta_i \ket{i} \colon \beta_i \ge 0, \ket{\phi}^{\otimes 2 k \ell}$)}
     \begin{itemize}
       \item If $\textup{SuperProduct}(\ket{\psi},(\ket{\phi}^{\otimes k})^{\otimes \ell})$ fails, then \emph{reject}.
       \item For $i=1,\ldots,\ell$ \begin{enumerate}
          \item Run the $\QMA^+(k)$ protocol $\calP$ on a new copy of $\ket{\phi}^{\otimes k}$.
          \item If protocol rejects, then \emph{reject}.
       \end{enumerate}
       \item \emph{Accept}.
     \end{itemize}
     \caption{Symmetric Simulator}\label{proc:sym_protocol}
\end{algorithm}

We now analyze the completeness and soundness of this simulation.

\begin{lemma}\label{lem:sym_simulator}
  Suppose $\calP$ is a $\QMA^+(k)$ protocol with completeness $c$ and soundness $s$.
  %
  Let $p(n)$ be a non-decreasing function such that $p(n) \ge C_0$ for a sufficiently large
  constant $C_0 > 0$. If $\ell \ge 8 p(n)^2 \ln(2)$ and $s \le 1/8p(n)^2$, then 
  SymSimulator has completeness $c^\ell$ and soundness at most $1/2 + 1/p(n)$.
\end{lemma}

\begin{proof}
  In the completeness case, we can assume that the proofs $\ket{\phi}$ have nonnegative amplitudes
  and $\ket{\psi} = \ket{\phi}^{\otimes k}$. Thus, SymSimulator accepts with probability at least $c^\ell$.

  Now, suppose that we are in the soundness case. Set $\epsilon =\abs{\braket{\psi}{\phi^{\otimes k}}}^2$.
  By~\cref{lem:super_product}, the super product test accepts with probability
    \[ 
  \left( \frac{ \epsilon\ell}{\ell+1} + \frac{1}{\ell+1} \right)\mper
    \] 
      Since $\ell \ge 2p(n)$, if $\epsilon < 1/2 + 1/2p(n)$, then the acceptance probability due to the super product test alone is at most $1/2+1/p(n)$
  and we are done. Therefore, from now on, we assume that $\epsilon \ge 1/2 + 1/2p(n)$. 
  
  Suppose $\ket{\phi} = \sum_i \alpha_i \ket{i}$, and let $\ket{\phi_+} = \sum_{i} \abs{\alpha_i} \ket{i}$. Thus, $\ket{\phi_+}$ is a valid nonnegative-amplitudes state.
  Since $\ket{\psi}$ has nonnegative amplitudes by assumption, we should have 
    \begin{equation}\label{eq:phi-plus-corr}
        |\langle \psi \mid \phi_+^{\otimes k} \rangle|^2 \ge |\langle \psi \mid \phi^{\otimes k} \rangle|^2 = \epsilon.
    \end{equation}
  This is because the latter inner product incurs some cancellations due to negative values, which are avoided in the former inner product. 
  (\ref{eq:phi-plus-corr}) together with \cref{fact:overlap-triangle-ineq} implies that  
    \begin{equation*}
        |\langle \phi^{\otimes k}, \phi_+^{\otimes k}\rangle|^2
            \ge 2\epsilon-1.
    \end{equation*}
  Since we are assuming $\epsilon > 1/2$, the trace distance between $\ket\phi ^{\otimes k}$ and $\ket{\phi_+}^{\otimes k}$ can be bounded as below
    \begin{equation}
        D(\phi^{\otimes k}, \phi_+^{\otimes k})\le 2\sqrt{\epsilon(1-\epsilon)}
    \end{equation}
  Note that $\calP$ accepts $|\phi_+^{\otimes k}\rangle$ with probability at most $s$ by the soundness of $\calP$.
  Therefore, each execution of the protocol $\calP$ on $\ket\phi ^{\otimes k}$ accepts with probability, by \cref{fact:trace_norm_acc}, at most 
  \[\min\set{1,2\sqrt{\epsilon(1-\epsilon)} + s}.\]
  %
  The overall soundness of SymSimulator becomes
    \[ 
  \left(\epsilon \frac{\ell}{\ell+1} + \frac{1}{\ell+1} \right) \left( \min\set{1,2\sqrt{\epsilon(1-\epsilon)} + s} \right)^{\ell} .
    \] 
 Now take  $\epsilon \ge 1/2 + 1/2p(n)$, and compute, we have
    \[ 
  2\sqrt{\epsilon(1-\epsilon)} \le 2 \sqrt{\frac{1}{4} - \frac{1}{4 p(n)^2}} \le 1 - \frac{1}{2 p(n)^2} + O\left(\frac{1}{p(n)^4}\right) \le 1 - \frac{1}{4p(n)^2}\mcom
    \] 
  where the last inequality relies on $p(n) \ge C_0$ for a large enough constant $C_0 > 0$.
  Using that $s \le 1/8p(n)^2$ and $\ell  \ge 8 p(n)^2 \ln(2)$, the final acceptance probability is
    \[ 
  \left(2\sqrt{\epsilon(1-\epsilon)} + s \right)^{\ell} \le \left(1 - \frac{1}{8p(n)^2} \right)^\ell \le \frac{1}{2},
    \] 
  concluding the proof.
\end{proof}

To remove  \ref{enu:many-general-states} the symmetric assumption of having multiple identical copies of $\ket{\phi}$ in SymSimulator, we use the PAPO channel $\Lambda$ and the PAPO channel takes
just two unentangled proofs $\ket{\phi'}$ and $\ket{\phi''}$ (of arbitrary amplitudes) as its input. In other words, we now simulate a $\QMA^+(k)$ protocol $\calP$ with:
\begin{description}
    \item \; (i)\, one nonnegative-amplitudes proof $\ket{\psi}$;
    \item (ii')\,  two general states $\ket{\phi'}, \ket{\phi''}$.
\end{description} 
A formal description of the new simulation is given in~\cref{proc:protocol}.

\begin{algorithm}[H]{Simulator}{($\calP$, $\ket{\psi} = \sum_i \beta_i \ket{i} \colon \beta_i \ge 0, \ket{\phi'}, \ket{\phi''}$)}
     \begin{itemize}
       \item Let $\rho$ be the output of our disentangler $\Lambda(\phi'\otimes \phi'')$ (i.e. \cref{thm:channel}).
       \item If $\textup{SymSimulator}(\calP, \ket{\psi},\rho)$ accepts, then \emph{accept}; else \emph{reject}.
     \end{itemize}
     \caption{Protocol Simulator}\label{proc:protocol}
\end{algorithm}

The analysis of~\cref{proc:protocol} is similar to that of \cref{lem:sym_simulator}. Therefore, instead of presenting an analysis of~\cref{proc:protocol} in isolation, we now apply this simulation for a $\QMA^+(k)$ protocol $\calP$ that solves a $\NEXP$-complete problem.  
In particular, we will need the following characterization of $\QMA^+(2)$ from~\cite{JW23} as shown in the following theorem.

\begin{theorem}[\cite{JW23}]\label{thm:qma2_plus_nexp}
  $\QMA^+(2) = \NEXP$.
\end{theorem}

\cref{proc:protocol} gives rise to a protocol for NEXP that improves the above theorem in two aspects. First, the new protocol uses three unentangled proofs among which only one is required to have nonnegative amplitudes. Second, the completeness and soundness gap of this protocol is about $1/2$. This seemingly mediocre gap is in fact a critical point, which we discuss in the next section.

\TheoQMAThree*

\begin{proof}
  From~\cref{thm:qma2_plus_nexp}, we apply the standard gap amplification by asking for more unentangled proofs to obtain a $\QMA^+(k)$ protocol $\calP$ with completeness $c = 1-\exp(-\poly(n))$
  and soundness $s=\exp(-\poly(n))$, where $k=\poly(n)$. Simulate $\calP$ using~\cref{proc:protocol}.
  By \cref{thm:channel}, $\rho = \Lambda(\phi'\otimes \phi'')$ is $1/\poly(n)$-close to a convex combination of
  product states $\int \ket{\phi} \bra{\phi}^{\otimes 2k\ell} d\mu$ with $\ell = \poly(n)$. Invoking the symmetric simulator, by~\cref{lem:sym_simulator},
  the completeness becomes $c^\ell \ge 1-\exp(-\poly(n))$ and the soundness $1/2+1/\poly(n)$ for a suitable choice of polynomial $\ell=\poly(n)$.
\end{proof}



\subsection{Almost-$\QMA^\R(k) = \NEXP$}
Next, we show how to go from the nonnegative amplitudes assumptions to almost general amplitudes.
Recall that the $\ell_2$-sign bias of a state $\ket{\psi} = \sqrt{a} \ket{\psi_+} + \sqrt{1-a}\ket{\psi_{-}}$,
where $\ket{\psi_+}$ and $\ket{\psi_{-}}$ are the normalized nonnegative and negative amplitudes parts of $\ket{\psi}$,
is defined as $\abs{a-(1-a)}$ (see~\cref{def:ell2_sign_bias}). 

\TheoAlmostGenAmp*

\begin{proof}
  We start with the $\QMA^+(3)$ protocol from~\cref{thm:qma3} with two general proofs $\ket{\phi'},\ket{\phi''}$
  and only one nonnegative proof $\ket{\psi}$. Let $M$ be the verifier measurement.
  In the completeness case, we can assume that $\ket{\psi}$ has nonnegative amplitudes so we proceed
  to analyze the soundness case.

  %
  In the \textup{almost}-$\QMA^\R(3)$ protocol, $\ket{\psi}$ will no-longer be assumed to have nonnegative
  amplitudes. Instead, we write  $\ket{\psi} = \sqrt{a} \ket{\psi_+} + \sqrt{1-a} \ket{\psi_{-}}$,
  where $\ket{\psi_+}$ and $\ket{\psi_{-}}$ are its nonnegative- and negative-amplitudes normalized states.
  Without loss of generality, suppose that $a \ge 1/2$.
  Furthermore, under the $\ell_2$-sign bias assumption, we may assume that 
  \begin{equation}\label{eq:a-bias}
      a \ge 1/2 + \sqrt{100/p(n)}.
  \end{equation}
  %
  Let $\ket{\phi'}$ and $\ket{\phi''}$ be some quantum states (ignoring the $\ell_2$-bias requirement) as to be used in the simulation~\cref{proc:protocol}.
  %
  The combined proofs of the \textup{almost}-$\QMA^\R(3)$ protocol can be expressed as $\ket{\xi} = \sqrt{a} \ket{\xi_0} + \sqrt{1-a} \ket{\xi_1}$, where
  $\ket{\xi_0} = \ket{\phi'} \otimes \ket{\phi''} \otimes \ket{\psi_+}$ and $\ket{\xi_1} = \ket{\phi'} \otimes \ket{\phi''} \otimes \ket{\psi_{-}}$. 
  %
  Denote s the soundness of $\QMA^+(3)$ protocol from~\cref{thm:qma3}. Then we can assume 
  \begin{equation}\label{eq:soundness-qma3}
  s \le 1/2 + 6/p(n).
  \end{equation}
  %
  Calculating the accepting probability of $M$ on $\xi$,
  \begin{align}
    \bra{\xi} M \ket{\xi} 
        &= a \bra{\xi_0} M \ket{\xi_0} + (1-a) \bra{\xi_1} M \ket{\xi_1} \nonumber \\
        %
        &\qquad\qquad +\sqrt{a(1-a)} \bra{\xi_0} M \ket{\xi_1} + 
        \sqrt{a(1-a)} \bra{\xi_1} M \ket{\xi_0}
        \nonumber \\
        %
        &\le s + \sqrt{a(1-a)} \left(\bra{\xi_0} M \ket{\xi_0} + \bra{\xi_1} M \ket{\xi_1}\right) 
        \nonumber\\
        %
        &\le (1+ 2\sqrt{a(1-a)})s. \label{eq:soundness-almost-qma}
  \end{align}
  where the first inequality follows from $M$ being PSD, \ie since $(\bra{\xi_0}-\bra{\xi_1}) M (\ket{\xi_0}-\ket{\xi_1}) \ge 0$
  implies $\bra{\xi_0} M \ket{\xi_0} + \bra{\xi_1} M \ket{\xi_1} \ge \bra{\xi_0} M \ket{\xi_1} + \bra{\xi_1} M \ket{\xi_0}$.
  By~(\ref{eq:a-bias}) and~(\ref{eq:soundness-qma3}), we have
    \[ 
  (1+ 2\sqrt{a(1-a)})s \le \left(2- \frac{8}{p(n)}\right) s \le\left(2- \frac{8}{p(n)}\right) \left(\frac{1}{2} + \frac{1}{p(n)}\right) \le 1 - \frac{2}{p(n)}.
    \] 
    

Note that by a suitable choice of polynomial $p(n)$ and  the initial completeness $c=1-\exp(-\poly(n))$ of the $\QMA^+(3)$ protocol of~\cref{thm:qma3},
we obtain a gap of $\Omega(1/p(n))$. To conclude the proof, we apply standard gap amplification using $k=\poly(p(n))$ proofs in $\textup{almost-QMA}^\R(k)$.
\end{proof}

We emphasize an important observation following from the above analysis: The ``half'' gap amplification in~\cref{thm:qma3} is almost optimal. A larger gap in~\cref{thm:qma3} by an additive term $1/\poly(n)$ (\eg if the soundness was at most $1/2-1/\poly(n)$) would allow us to completely discard the $\ell_2$-sign bias assumption in \cref{thm:almost_gen_amp}, showing $\NEXP=\QMA^{\mathbb{R}}(k)$. This can be easily seen in (\ref{eq:soundness-almost-qma}), when $s<1/2-1/\poly(n)$, the RHS will be at most $1-1/\poly(n).$ It means that $s=1/2\pm 1/\poly(n)$ in~\cref{thm:qma3} is a critical point. In the case that $\QMA(k)^\R\not=\NEXP$, there is a sharp phase transition.
